\section{Rank Tracks $K$ and Is Causally Necessary}
\label{sec:recruit}

\textbf{Recruitment.} Without any rank term in the loss, the trained
$d{=}16$ state's effective rank rises one-for-one
with binding count across a ten-point grid from $K{=}1$ to $K{=}16$
(Spearman $\rho = 1.0$; 2.42 at $K{=}1$, 8.20 at $K{=}8$, 15.09 at
$K{=}16$), leaving the pre-registered $[0.7K, 1.3K]$ band only at
$K{=}1, 2$, by overshoot; beyond $K = d$ rank saturates and declines.
% <!-- evidence: R1 -->
A subsequent, larger replication sweep is directionally identical.
% <!-- evidence: R1b -->

\begin{figure}[t]
\centering
\includegraphics[width=0.92\textwidth]{fig_forcerank.pdf}
\caption{The primary, pre-registered causal test: $\recninety$ under
train-time force-rank $k$, regenerated from the archived sweep (dashed
line at $k = K$). The step lands at $k \approx K$ in every cell whose
unconstrained model reaches exact recovery.}
% <!-- evidence: R2 -->
\label{fig:forcerank}
\end{figure}

\textbf{Causal necessity.} Training with $Z$ constrained to rank $k$
throughout (Figure~\ref{fig:forcerank}): at $d{=}8, K{=}4$
the step is exact, $k \leq 3$ yields $\recninety \leq 0.0004$ and
$k = 4$ yields 0.97, a discontinuity at the provable bound.
% <!-- evidence: R2 -->
The $d{=}16$ cells show the same knee near $k \approx K$ with mild
post-knee non-monotonicity from seed noise; the replication resolves
the $K{=}12$ transition into a ragged ramp, not a clean step.
% <!-- evidence: R2b -->
The two halves answer \citet{larson2026gradient}: the earlier
rank-blindness was a property of a rank-1-solvable task, not of the
gradient.
